\chapter*{APÉNDICE}
\addcontentsline{toc}{chapter}{LECTURA I - El timbre y el carácter del violín}
En toda vida, las explicaciones de términos técnicos resultan útiles en algún momento; y este momento puede llegar durante la juventud o posponerse hasta la vejez. Me resulta imposible escribir extensamente sobre las peculiaridades del sonido del violín sin emplear algunos términos técnicos; y, recordando la época en que el lenguaje técnico me resultaba confuso, el placer que sentía al explicarlo y mi deseo de que todos los lectores reciban una recompensa por el esfuerzo que supone leer los detalles áridos de las páginas anteriores, por lo tanto, incluyo con gusto una explicación de algunos de los términos más confusos que se emplean aquí. Para aquellos lectores que no estén interesados en dichas explicaciones, no es necesario sugerirles que pierdan el tiempo leyendo estas últimas páginas.
ALT: personal.	Todos los tonos en la primera octava por encima de la
ALTISSIMO: Todos los tonos por encima de alt.
CONDICIONES METEÓRICAS: Se refiere a condiciones atmosféricas como la presión barométrica, la temperatura, la humedad (vapor de agua), los vientos y las nubes. Dado que todos estos factores, individualmente o en conjunto, influyen en la distancia que recorren las ondas sonoras, resultan de gran interés para el estudiante de la sonoridad del violín.
PRESIÓN BAROMÉTRICA: Se refiere al peso del aire indicado por la acción del mercurio en el tubo del barómetro. A medida que el mercurio desciende, el peso del aire es menor. A medida que el mercurio asciende,
289

El peso del aire es mayor. Estas variaciones en el peso del aire son frecuentes, especialmente durante los meses de verano, y pueden ser tan grandes que provocan grandes diferencias en la distancia que recorre cualquier sonido, independientemente de su origen. Curiosamente, el violín, de entre todos los instrumentos musicales, es el más sensible a las condiciones atmosféricas, como se manifiesta en la debilidad de la intensidad del sonido en fechas con exceso de calor y vapor de agua. En tales condiciones, el timbre de las cuerdas se degrada inevitablemente, sin importar quién sea el fabricante ni la calidad del material.
LA NUBE NIMBUS: Extendiendo todo
El cielo, funciona para aumentar la distancia que recorren las ondas sonoras.
TEMPERATURA: Puede ser muy alta o muy baja.
lo suficientemente bajo como para disminuir considerablemente la distancia recorrida por las ondas sonoras.
VIENTOS: Operan para disminuir la distancia recorrida por las ondas sonoras en oposición a la corriente y para aumentar la distancia recorrida con la corriente.
actual.
HORA DEL DÍA: Es importante para registrar la distancia que recorren las ondas sonoras, porque durante las horas del mediodía el sonido se propaga con mayor dificultad que en cualquier otro día.
AGENTES REFLEJANTES DEL SONIDO: ¿Existen objetos como edificios, colinas o madera que contribuyen a aumentar la distancia que recorre el sonido del violín en la prueba de intensidad a larga distancia al aire libre?
Por lo tanto, el registro de cualquier violín en cuanto a "potencia de transmisión" posee poco interés a menos que esté acompañado de...
290

Se determina mediante las lecturas del barómetro, el higrómetro y el termómetro, junto con la dirección y la velocidad del viento, el grado de nubosidad, la hora del día y la presencia o ausencia de agentes reflectantes del sonido. Con este registro completo, se puede garantizar con seguridad que un violín repetirá su interpretación en condiciones meteorológicas similares. Es importante para el violinista conocer la distancia a la que el sonido de su violín viaja con facilidad. De esta manera, el intérprete puede evitar el esfuerzo excesivo con el arco, ya que dicho esfuerzo siempre produce un efecto desagradable en el oído musicalmente cultivado.
OSCILACIÓN, AMPLITUD DE OSCILACIÓN, VIBRACIÓN, tanto NORMAL como TANGENCIAL, NODOS, SEGMENTOS VENTRALES, SOBRETONOS ARMÓNICOS y SOBRETONOS DISONANTES: Se explican así: Limitando este asunto a hechos de valor práctico para el sonido del violín, empleo estas dos cuerdas como medio para ayudar a aclarar las definiciones a la comprensión de cada lector. Estas cuerdas tienen la misma longitud, el mismo diámetro y son del mismo material, pero son desiguales en perfección estructural. La cuerda A es aparentemente perfecta, la cuerda B es aparentemente imperfecta. La acción de estas cuerdas explicará no solo algunos términos técnicos aplicados a los cuerpos vibrantes, sino también la necesidad de emplear madera de perfección estructural para producir el violín de los mejores o más altos valores tonales; también explicará por qué el luthier, científico o no, debe enfrentarse a la derrota cuando trabaja con madera de imperfecciones estructurales. 291

Se dice que un cuerpo inmóvil está en reposo. Cuando un cuerpo se desplaza una cierta distancia desde el reposo, regresa a él, recorre una distancia igual en dirección opuesta y vuelve por segunda vez al reposo, se dice que dicho cuerpo vibra u oscila, según se prefiera. La distancia recorrida por el cuerpo desde su reposo es la amplitud de su oscilación.
[Resulta interesante observar que los filósofos ingleses y alemanes difieren de los franceses en cuanto a qué movimiento de un cuerpo vibrante completa una vibración. Los primeros sostienen que una vibración completa consiste en un movimiento completo en cada sentido desde el punto de reposo; mientras que los segundos sostienen que un movimiento desde el punto de reposo con un retorno completa una vibración. Así, según el método de expresión francés, la afinación de concierto se da como La igual a 900 vibraciones por segundo; mientras que, según los métodos inglés y alemán, la afinación de concierto se da como La igual a 450 vibraciones por segundo.]
Al fijar estas cuerdas por separado a bloques inamovibles y a clavijas separadas, se aplica una tensión igual. La aplicación de un arco de violín hace que la cuerda A se enrolle rápidamente alrededor de su eje longitudinal; y, dicho enrollamiento continúa en una dirección hasta que la energía elástica en la fibra de la cuerda supera la fricción del arco; entonces, la cuerda se desenrolla, solo para volver a enrollarse instantáneamente. Este rápido enrollamiento y desenrollamiento produce fuertes golpes sobre las cuerdas contiguas.
moléculas de aire, y, como tales moléculas son elásticas
292

esferas, y, como tales esferas se tocan entre sí, por lo tanto, los golpes de la cuerda provocan un movimiento de ondas sonoras; y; como tal movimiento llega a nuestros tímpanos, (tímpanos, ) nos damos cuenta de un tono que procede de la cuerda; y, como tal tono es producido por la acción de toda la longitud de la cuerda, por lo tanto este tono es del tono más bajo posible para una cuerda de esta longitud, y, con esta tensión; por lo tanto este tono es el tono fundamental de la cuerda A.
[Aunque este tono es fundamental para la cuerda A, no podemos suponer que no se pueda producir ningún otro tono fundamental en la cuerda A, porque, al acortar la cuerda "deteniéndola", como con la presión del dedo, se produce un nuevo tono de tono más alto; y dicho nuevo tono puede llamarse legítimamente el tono fundamental de una nueva tonalidad; y lo mismo se aplica a todos los demás tonos de
tono más agudo.]
La acción de la cuerda A provoca de inmediato la aparición de fenómenos absorbentes; uno de los cuales consiste en que la cuerda describe un círculo alrededor de su eje longitudinal. Observamos que dicho círculo es más pequeño en los extremos de la cuerda y más grande en el punto medio. Sin ningún razonamiento, es evidente que la mayor fuerza en los golpes de esta cuerda se aplica en el punto de su mayor amplitud de oscilación; y que la menor fuerza
Sus golpes se producen en los puntos de menor amplitud de oscilación. Asimismo, es evidente que el punto de mayor amplitud de oscilación no puede darse en ningún otro lugar que no sea el punto medio.
293

[Actuando sobre este hecho, al graduar la caja de resonancia para obtener la máxima potencia tonal, el grosor, desde la posición del puente y desde los extremos de la placa, disminuye ligeramente hasta el punto medio; y, con el propósito de aumentar la amplitud de oscilación en dicho punto. Por la acción de la cuerda, es claro que no se puede asegurar la misma amplitud de oscilación en ningún otro punto entre la posición del puente y los extremos de la placa.
Otro fenómeno, observado en la cuerda A, dirige la atención a los nodos, segmentos ventrales y armónicos. Al detener ligeramente la cuerda a la mitad, en cualquier dirección se observan varios puntos donde la cuerda está en reposo. Estos puntos de reposo se llaman nodos. Están igualmente distantes entre sí, y el espacio entre dos nodos se llama segmento ventral; y observamos que cada segmento ventral actúa precisamente como actúa toda la longitud de la cuerda; es decir, la amplitud más amplia de oscilación del segmento ventral está en su punto medio; y que la acción de estos segmentos ventrales produce un tono musical. Como los segmentos ventrales tienen una longitud similar, el tono de cada segmento es del mismo tono que el tono de su vecino; y como el tono de dicho tono está en armonía con el tono fundamental, dicho tono se llama armónico, y al unirse con el tono fundamental, produce lo que se llama un tono "rico", una cualidad de tono rara, una cualidad de tono muy valorada por el violinista solista, una cualidad de tono en la que no
294

otro dispositivo musical se acerca al "mejor" violín una calidad de tono imposible de lograr con la caja de resonancia de imperfecciones estructurales; y la razón se demostrará mediante la acción de la estructura.
cuerda B imperfecta.
Aplicados a la música, los términos «consonante» y «disonante» varían en significado tanto como «santo» y «diablo» aplicados a la humanidad. Los armónicos consonantes y los armónicos graves son la causa de la encantadora belleza del sonido del violín. Se sitúan entre términos como «dulzura» y «riqueza». A ellos les debemos la corona que con tanto entusiasmo ofrecemos al «Rey». Pero se nos recuerda, muy a menudo, que existen otras coronas. Satanás lleva una. Curiosamente, la corona de Satanás podría provenir del violín; y aún más extraño, la actividad de Satanás en acaparar material para su corona se ve complementada por fabricantes de violines tan científicos que ignoran la perfección estructural de la madera de la caja de resonancia.
Al aplicar el arco a la cuerda B, el armónico disonante (joya principal en la corona de Satanás) se destaca claramente para la percepción del sentido del oído; y la causa se destaca claramente para la percepción del sentido de la vista. Observamos que los nodos en este agente productor de tono estructuralmente imperfecto existen a intervalos variables; por lo tanto, los segmentos ventrales tienen longitudes variables; por lo tanto, los tonos de estos segmentos ventrales tienen un tono aleatorio; por lo tanto, muchos de esos tonos no están en armonía con el tono fundamental de la cuerda; por lo tanto, el sonido de este agente productor de tono estructuralmente imperfecto...
295
\label{tono_frio} \label{sobretonos_dis2}
El agente productor de tono de los pies es RUIDO; y, por la misma razón, el tono de la caja de resonancia estructuralmente imperfecta es RUIDO; por lo tanto, las letras en la corona de Satanás son N-O-ISE.
En los libros de texto dedicados a la filosofía del sonido musical, la caja de resonancia del violín se compara con un "haz de cuerdas"; y, determinando por el tono "rico" de la cuerda perfecta, y por el tono "rico" de la caja de resonancia perfecta, se aclara que parte de la acción vibratoria de la cuerda perfecta se duplica en la fibra de la caja de resonancia perfecta; también, que parte de la acción de la cuerda imperfecta se duplica en la acción de la fibra de la caja de resonancia imperfecta. Así, en la caja de resonancia de fibra perfecta, los nodos aparecerán a intervalos regulares; por lo tanto, los segmentos ventrales tendrán longitudes uniformes; por lo tanto, los tonos de los segmentos ventrales tendrán un tono uniforme; por lo tanto, dichos tonos están en armonía con el tono fundamental; de ahí el tono "rico" del violín. Así, en la caja de resonancia de fibra imperfecta, como nudos, rizos, densidad desigual, desviación de la fibra de las líneas rectas, gran densidad e inelasticidad del tejido conectivo, albura y manchas negras, los nodos deben aparecer a intervalos irregulares; Por lo tanto, los segmentos ventrales deben tener longitudes variables; por lo tanto, los tonos de los segmentos ventrales deben tener un tono aleatorio; por lo tanto, muchos de esos tonos deben estar fuera de armonía con el tono fundamental; de ahí el violín ruidoso que desafía a la ciencia; el violín que desafía al luthier más hábil; y, cuando se suprime el ruido en él,
296

es el violín de tono "frío"; y el violín de
El tono "frío" ya no es "El Rey".
MOVIMIENTO VIBRATORIO: En el violín, la caja de resonancia despierta un profundo interés en el estudiante; y este interés se debe a que, junto con otros factores, el timbre del violín depende tanto del movimiento vibratorio normal, que se propaga a lo largo de las fibras, como del movimiento vibratorio tangencial, que se propaga a través de ellas. Estos dos movimientos vibratorios son la base misma del timbre del violín. Son los factores que enriquecen el sonido de las cuerdas. Sin estos factores, el timbre del violín se degrada hasta el nivel de un violín de palo de escoba. Al colocar arena seca sobre la caja de resonancia plana, se definen tanto el movimiento vibratorio normal como el tangencial. Distribuida uniformemente sobre dicha caja de resonancia, la arena es impulsada hacia arriba por la fuerza de oscilación de la placa, alcanzando una altura proporcional a la amplitud de dicha oscilación. Es aquí donde observamos variaciones asombrosas en la elasticidad de diferentes muestras de madera. Aquí también se muestra la gran diferencia entre la acción de la fibra estructuralmente perfecta e imperfecta. Así: Sobre la caja de resonancia perfecta, su oscilación, continuada durante un período de tiempo que varía según la acción del resorte, obliga a la arena a abandonar los segmentos ventrales y a acumularse en los nodos. De este modo, se determina la regularidad en la distancia entre los nodos. Así, se muestra la irregularidad en la distancia entre los nodos en la fibra estructuralmente imperfecta.
297

caja de resonancia.
Es cierto que esta demostración no puede realizarse con éxito en la caja de resonancia cóncavo-convexa del violín. También es cierto que dicha demostración es innecesaria. Es evidente que estos hechos son similares en estas dos formas de caja de resonancia. En ambos casos, la presencia de las salidas limita el recorrido vibratorio transversal, de la siguiente manera: Cuando los golpes de las cuerdas se aplican sobre la caja de resonancia en la posición del puente, el movimiento vibratorio en la caja de resonancia comienza en dicha posición; y, debido a la proximidad de las salidas, es evidente que el movimiento vibratorio por encima del puente se limita a la vibración normal hasta después de pasar el extremo superior de las salidas. Debajo del puente, en el lado de los graves, la distancia a la salida es suficiente para permitir la acción simultánea de la vibración normal y transversal. Debajo del puente, en el lado de los agudos, el poste detiene eficazmente el recorrido de ambos movimientos vibratorios, como se demuestra satisfactoriamente dividiendo el cuarto inferior derecho de la caja de resonancia. Basándonos en las conclusiones extraídas de la evidencia aportada por cientos de cajas de resonancia usadas, resulta evidente que el tono "rico" depende de la propagación sin obstáculos de la vibración normal; asimismo, que el volumen del tono depende en gran medida de la vibración transversal.
Un hecho interesante se muestra en el registro de velocidades comparativas adjuntas a estos movimientos vibratorios. Así: En el pino, vibraciones normales,
3322 metros
por segundo. En pino, transversal
vibraciones, 1405 metros por segundo. 298	Por lo tanto,

Mientras que la vibración transversal se propaga una unidad, la vibración normal se propaga entre 2,36 y 100 unidades; sin embargo, la evidencia obtenida de cajas de resonancia usadas indica que esta proporción no es constante, sino que se ve muy modificada por las peculiaridades estructurales de la veta en diferentes muestras de madera. Por lo tanto, la vibración transversal se ve atenuada tanto por la veta extremadamente blanda como por la extremadamente densa.
En conjunto, estos movimientos vibratorios constituyen la base tanto del volumen como de la calidad del sonido. En la producción del volumen del sonido a partir de la caja de resonancia, la vibración transversal ejerce la influencia más dominante. Esto se demuestra fácilmente de la siguiente manera: con una caja de resonancia demasiado rígida para la fuerza de las cuerdas de cierto calibre, como el calibre 2, la disminución gradual del grosor desde la unión central hacia los bordes aumenta la distancia recorrida por la vibración transversal; por lo tanto, aumenta el área de la superficie de golpeo; en consecuencia, un mayor número de moléculas de aire dentro del cuerpo del violín reciben el mismo golpe; por lo tanto, se obtiene un mayor volumen de sonido.
[Es fácil aumentar el volumen del sonido del violín hasta el punto de arruinar la intensidad del tono.]
La pérdida de potencia tras su uso es un tema de gran interés para el estudioso de los fenómenos tonales del violín. He observado que dicha pérdida se debe a los siguientes factores:
1. Aumento de la rugosidad de las superficies interiores no protegidas.
2. Desintegración sobre una superficie desprotegida.
3. Aumentar la flexibilidad del foro de debate
299

uso.
4. Grado de densidad de fibras y resistencia natural del tejido conectivo.
5. Cantidad de uso y vigor del arco.
Aquí radica la importancia de la flexibilidad. Debido a su rigidez inherente, las partes duras de la fibra de la caja de resonancia no pueden doblarse transversalmente; sin embargo, el tejido conectivo, al ser elástico tanto en longitud como en anchura, permite la flexión transversal. Por lo tanto, la vibración transversal depende de la elasticidad del tejido conectivo, la cual disminuye con el uso. La tasa de disminución se ve modificada por el grado de flexión, la frecuencia de flexión y la resistencia natural del tejido conectivo, característica de cada tipo de madera. La resistencia del tejido conectivo en la madera de la caja de resonancia del violín varía considerablemente; y, por esta razón, la durabilidad del sonido del violín presenta periodos muy variables. Estos hechos conocidos pueden aplicarse a la caja de resonancia del violín de la siguiente manera: una caja de resonancia que solo utiliza una parte de su fuerza elástica durará más que una que utiliza el límite máximo de dicha fuerza.
Ambas tapas del violín pueden actuar como resortes. Esto depende del grosor y la rigidez inherente. Evidentemente, un mayor volumen del sonido del violín depende de una reducción en la rigidez de la tapa que permita una mayor distancia recorrida y una mayor amplitud de la vibración transversal; asimismo, es evidente que esta mayor amplitud se logra mediante una mayor flexión de la tapa, lo que acelera la pérdida de la fuerza elástica.
300

Estos hechos señalan inequívocamente el riesgo de otorgar un volumen de sonido marcado a un violín nuevo, ya que este es propenso a sufrir una pérdida perjudicial de potencia sonora a una edad que bien podría considerarse prematura. Desde mi punto de vista, es mucho más prudente limitar ligeramente tanto la distancia como la amplitud de la vibración transversal y confiar en que la flexibilidad de la madera aumente inevitablemente con el uso.
[La obra del profesor Pietro Blaserna, de la Real Universidad de Roma, es un admirable tratado para quienes deseen estudiar el sonido en relación con la música.]
POSICIÓN DEL POST: Como el problema del post ha recibido diversas interpretaciones por
Dado que existen diferentes filósofos y no se vislumbra una base autorizada sobre la cual fundamentarse, todas las personas son libres de sostener y expresar opiniones sobre esta cuestión según les convenga. Sobre esta base, se presentan los siguientes cálculos: Como la fuerza de las cuerdas del violín depende del aumento producido por la vibración normal y transversal en la caja de resonancia, más la acción simpática debida a la proximidad de las cuerdas, y como la colocación del poste debajo del puente opera para confinar dichas actividades al cuadrante superior derecho, como se demuestra al dividir el cuadrante inferior derecho, por lo tanto, la posición del poste, en relación con el puente, rige la ubicación de las actividades de la caja de resonancia.
aumentando la fuerza de las cuerdas A y E; y, como la colocación del poste sobre el puente opera para transferir la vibración normal y transversal al cuarto inferior derecho de la caja de resonancia, y, como dist-
301

La anulación anula la acción simpática, por lo tanto, colocar el poste sobre el puente funciona para aplicar golpes de fuerza reducida sobre el aire contenido;
Por lo tanto, se produce una disminución de la potencia de los tonos A y E.
Dado que la influencia del poste recae principalmente sobre las cuerdas A y E, al seleccionar su posición deben tenerse en cuenta ciertos efectos, como la máxima potencia y la mejor calidad de sonido. Sucede con frecuencia que la máxima potencia perjudica la calidad del sonido; por lo tanto, en tal caso, debe determinarse un equilibrio entre ambas. Para ello, es necesario considerar el poste en sí, independientemente de su posición.
Esta consideración es necesaria porque cualidades físicas del mástil, como la longitud, la densidad y la masa, influyen notablemente en los tonos La y Mi. Cualquiera de estos atributos puede frustrar la intención. Evidentemente, la formación preparatoria resulta beneficiosa, de hecho, indispensable para obtener los mejores resultados en el ajuste del mástil. Como en todo lo relacionado con la modificación del timbre del violín, el sentido musical del ajustador debe tenerse en cuenta; y, cuando este sentido es deficiente, los fracasos en los intentos de ajuste no deben atribuirse a Stradivarius.
Además de las cualidades físicas del puesto, su posición se ve modificada por cada uno de los siguientes factores:
1. Altura del puente.
2. Altura del arco.
302

3. Espesor de las placas en la posición del puente.
4. Método de graduación.
5. Calidad elástica de las placas.
6. Diámetro y calidad de la cuerda.
Obviamente, la posición que garantice la mayor amplificación de los tonos A y E solo puede determinarse mediante prueba y error.
A partir de diferentes métodos para configurar la publicación, se presenta lo siguiente:
Partiendo de la premisa de que la longitud de las placas es de 14 pulgadas, el puente se coloca primero a 8 pulgadas del extremo superior de la caja de resonancia por las razones expuestas anteriormente.
[Esta posición del puente se modifica según el método de graduación.]
Aplicando suficiente tensión a las cuerdas para mantener el puente en posición, coloque el poste directamente debajo del pie derecho del puente y tense las cuerdas hasta el diapasón normal (tono internacional) o hasta el tono de concierto, manteniendo a partir de entonces el tono y el calibre de las cuerdas acordados.
Con el poste en esta posición, observe la potencia de los tonos La y Mi; luego, baje el poste 1/16 de pulgada y observe el aumento de potencia de La y Mi, y continúe así hasta llegar al punto donde la potencia de La y Mi parezca disminuir. Al regresar el poste al punto donde la potencia de La y Mi aumenta al máximo, muévalo 1/32 de pulgada.
Al llegar a ese punto, observe el equilibrio de poder entre A y E. Si E posee mayor poder, mueva el poste hacia la izquierda 1 pulgada y 32, y observe el resultado; continuando así
303

hasta que se establezca la igualdad de poder. Si el A
poseer el mayor poder, mover el puesto al
bien.
En este trabajo, la derrota puede derivarse del siguiente hecho: Con un diámetro del poste de 3/16 pulgadas, una ligera desviación de la perpendicular hace que la caja de resonancia descanse sobre uno u otro borde del poste. Así, cuando la caja de resonancia descansa sobre el borde inferior del poste, la distancia desde el puente es prácticamente 3/16 mayor de lo que parece; y como 1/32, más o menos, causa un efecto perceptible en el tono de las cuerdas A y E, por lo tanto 6/32 resulta suficiente.
Motivo de la derrota.
ÁNGULOS de INCIDENCIA y REFLEXIÓN,
\label{angulos}
Dentro del cuerpo del violín, existen factores importantes que influyen en la intensidad del sonido. La influencia de estos ángulos es fácil de comprender; sin embargo, estas líneas no se han trazado con precisión en ningún tipo de curvatura.
Estos ángulos dependen para su existencia de un cuerpo en movimiento y una superficie reflectante sólida, de la siguiente manera: La pelota, lanzada contra la pared, viaja a lo largo de la línea de incidencia; y, al rebotar, viaja a lo largo de la línea de reflexión. Si la línea de incidencia es perpendicular a la pared, entonces la pelota rebota directamente a su punto de partida; por lo tanto, las líneas de incidencia y reflexión son una; y así son estas líneas en el violín de placas perfectamente planas. Si ahora, el lanzador se coloca a un lado u otro de la línea que es perpendicular a la pared, y planta la pelota en un punto donde la perpendicular...
304

Si la línea ular toca la pared, entonces la bola no rebotará en la línea de incidencia ni en la línea perpendicular, sino que viajará a lo largo de una nueva línea que tiene exactamente el mismo ángulo con la pared que la línea de incidencia; y así son las líneas de viaje de las ondas sonoras dentro del violín de placas arqueadas. Es evidente que tales paredes interiores que desvían las líneas de viaje de las ondas sonoras lejos de las salidas causan pérdida de intensidad del tono; y por el contrario, tales paredes interiores que dirigen las líneas de viaje de las ondas sonoras
El viaje hacia las salidas provoca un aumento de la intensidad del tono proporcional al grado de concentración. Verbum savia Primero, establezca las paredes interiores;
Más tarde, las paredes exteriores.
ETIQUETA, BARNIZ y PRECIO vs TONO DULCE: Johann Holtzhammer es un pastor de ovejas. Su casa de verano está en la montaña. Desde la mañana temprano hasta el atardecer cubierto de rocío, Johann escucha dulces sonidos. El canto de los pájaros es dulce. El canto del arroyo ondulante es dulce. El susurro de las ramas que se mecen es dulce. El olor a pino es dulce. Incluso el ruidoso zumbido de las ruedas de la industria está desprovisto de la onda sonora antes de que llegue a los oídos de Johann. En la base de la roca desprendida,
Las bayas silvestres ofrecen su dulce apariencia y sus jugos aún más dulces. El aire de la montaña es dulce. El agua del arroyo es dulce. De hecho, para Johann la vida es un ciclo de dulzura; sin embargo, Johann anhela algo que mantenga sus manos ocupadas. Con la NATURALEZA palpitando a su alrededor, la ociosidad se vuelve opresiva. ¿Qué haría? ¿Qué podría hacer? ¿Hacer un violín? ¿Por qué no? Aquí está
305

Un pino gigante caído. Allá hay un arce muerto. Ahí está el hacha. Aquí hay una navaja. Donde hay voluntad, hay un camino. Esos troncos gigantes se han estado secando desde que el abuelo de Johann pastoreaba las ovejas; pero el hacha está afilada y la navaja de Johann siempre está lista para la acción. El tiempo no importa. El resultado lo es todo. Mientras Johann trabaja, una alondra cercana canta con indudable aprobación. Es cierto que Johann no sabe nada de la escala pitagórica; nada de la escala de Palestrina; nada del piano, ni de la escala temperada; nada de armónicos consonantes; nada de armónicos bajos; pero sí conoce el dulce sonido. Eso es todo lo que...
necesidades.
Johann mancha su violín con el jugo de las bayas de la montaña.
¡Suficiente!
Se desafía la imitación. El secreto muere con Johann.
El tono emocionante, encantador, conmovedor y dulcemente tierno del violín de Johann había sido durante mucho tiempo motivo de asombro y veneración por parte de los pájaros antes de su muerte; y, después de su muerte, cada pájaro de esa montaña competía con su vecino en cantar dulce y bajo en la fecha recurrente en que Johann partía con su violín. Ningún pájaro de esa montaña mencionó jamás los ángulos afilados del violín de Johann; tampoco lo hizo Caronte cuando Johann solicitó transporte a través del Estigia; tampoco lo hizo el guardián de la gran puerta, como aparece en la continuación. Cuando Johann entró en la larga avenida, una multitud de sombras con los ojos caídos estaban de pie en un
306

deteniéndose a cada lado. Algunos de ellos llevaban "obras maestras" en las que había etiquetas llamativas con las palabras ETIQUETA y VARIEDAD, y un número limitado llevaba cifras llamativas en la etiqueta, como 10.000 dólares y 12.000 dólares.
Aunque aún no era asunto de Johann, era un misterio por qué tanta demora. Como este portero solo necesita una milmillonésima parte de un instante para calcular el valor ad valorem en cualquier factura, en el instante en que llegó el envío inmortal de Johann, la gran puerta se abrió de golpe; y, en ese breve lapso, un destello de gloria recorrió la larga avenida; y, cuando la gran puerta se cerró tras Johann, se oyó desde la larga avenida un sonido como el crujir de dientes.
